\documentclass[11pt,a4paper]{article}
\usepackage[T1]{fontenc}
\usepackage[utf8]{inputenc}
\usepackage{amsmath,amssymb,amsthm,newpxtext,newpxmath,microtype}
\usepackage[a4paper,left=27mm,right=27mm,top=25mm,bottom=27mm]{geometry}
\usepackage{booktabs,array,longtable}
\usepackage[hyphens]{url}
\urlstyle{same}
\usepackage[colorlinks=true,linkcolor=black,urlcolor=black,citecolor=black]{hyperref}
\usepackage{fancyhdr}
\pagestyle{fancy}\fancyhf{}\fancyhead[L]{\small Erd\H{o}s \#1041: round-six research return}\fancyhead[R]{\small\thepage}
\renewcommand{\headrulewidth}{0pt}
\setlength{\headheight}{14pt}
\setlength{\parskip}{.22em}
\linespread{1.025}
\newtheorem{theorem}{Theorem}[section]
\newtheorem{proposition}[theorem]{Proposition}
\newtheorem{lemma}[theorem]{Lemma}
\newtheorem{corollary}[theorem]{Corollary}
\theoremstyle{remark}\newtheorem{remark}[theorem]{Remark}
\newcommand{\D}{\mathbb D}
\newcommand{\C}{\mathbb C}
\newcommand{\R}{\mathbb R}
\newcommand{\length}{\operatorname{length}}
\newcommand{\Area}{\operatorname{Area}}
\newcommand{\Rea}{\operatorname{Re}}
\newcommand{\K}{\mathcal K}
\newcommand{\clip}[1]{\left[#1\right]_{-1}^{1}}
\title{Scale-free bounds, quadratic coordinates, and metric cusps in polynomial lemniscates}
\author{Round-six mathematical and editorial return}
\date{7 September 2026}
\hypersetup{pdftitle={Erdos 1041: round-six research return},pdfauthor={Research return prepared with ChatGPT},pdflang={en-GB}}
\begin{document}
\maketitle
\begin{abstract}
The registered $13/25$ theorem already implies a degree-uniform connection
bound $(5/2)\mu^{1/n}$ in the smaller open level $(25/13)\mu$.
We extract that consequence and give an explicit, nonasymptotic connection
for the Blaschke-power polynomials
$[z(z+b)]^N-(1+bz)^N$, for every $N\ge3$ and $0<b\le1/(10N)$.
Inverting the quadratic factor preserves cancellation along two radial
pieces; an amplitude deficit pays for the intervening circular arc.
The registered degree-24 spectral witness consequently has a connection
shorter than $1/2$ at level at most $3750/3763$.
Weighted radial barriers also give a two-sided cusp law of order
$\delta^{1/(n-1)}$ for the optimal all-curve distance at $z^n-1$, along this
family in every fixed even degree $n\ge6$.
We distinguish these ordinary proofs from the supplied algebraic Lean
components, and give applicable revisions and a map of the remaining
geometric obstruction.
\end{abstract}

\noindent\textit{Scope.} The scaling result below is already a corollary of
the live corpus theorem. The sector construction, finite metric consequence
and two-sided cusp estimate are new ordinary arguments in this return.
They have not been independently reviewed or formalised in Lean, and no
novelty claim is established. Erd\H{o}s \#1041 is not claimed solved.

\section{The stronger general bound already in the live theorem}

For a squarefree monic polynomial $f$ of degree $n\ge2$, write
\[
 \mu(f)=\min_{f'(c)=0}|f(c)|>0,\qquad K_t(f)=\{z:|f(z)|\le t\}.
\]
The live Theorem \texttt{res:low-critical-thirteen-twentyfifths} states that
$\mu(f)\le13/25$ gives a curve shorter than two between distinct roots in
$\{|f|<1\}$. There is no root-location hypothesis. The ordinary source
\path{AngularBudgetLowCriticalClosure.md} already states its scaling form
as Corollary B$''$; the exact-result entry
\path{erdos1041_all_degree_mu_13_25_2026_09_05} includes the same conclusion.

\begin{proposition}[registered scaling consequence]\label{thm:scale}
Every squarefree monic polynomial of degree $n\ge2$ has two distinct roots
joined in $\{|f|<(25/13)\mu\}$ by a curve of length less than
\[
 2\bigl((25/13)\mu\bigr)^{1/n}.
\]
In every degree the length can be chosen less than $(5/2)\mu^{1/n}$.
\end{proposition}
\begin{proof}
Set $s=((25/13)\mu)^{1/n}$ and $g(z)=s^{-n}f(sz)$. Then $g$ is monic,
its critical points are $c/s$, and $\mu(g)=s^{-n}\mu=13/25$.
The live theorem gives a curve of length less than two in $\{|g|<1\}$.
Multiplication by $s$ proves the first claim.
For $n\ge3$,
\[
 (25/13)^{1/n}\le(25/13)^{1/3}<5/4,
 \qquad \frac{25}{13}<\frac{125}{64}.
\]
For $n=2$, write $f(z)=(z-a)^2-d^2$ with $d\ne0$.
The segment $a+td$, $-1\le t\le1$, has length $2|d|=2\sqrt\mu$
and satisfies $|f|\le\mu$.
\end{proof}

The sharper common coefficient supplied by this argument is
$2(25/13)^{1/3}=2.487113\ldots$; the rational $5/2$ is convenient for the
introduction. Neither is asserted optimal. If $f$ has a repeated root,
$\mu=0$ and a constant path in $K_0$ joins two listed occurrences. The strict
squarefree displays must not be extended to claim length less than zero or
containment in $\{|f|<0\}$.

This changes the paper's hierarchy. The $71/10$ theorem remains a valid
independent area argument and should be retained, as requested. It is no
longer the strongest threshold-free numerical bound visible from the paper's
own hypotheses. It uses the larger level $2\mu$ and the larger length
coefficient. Its expository value is the complete, transparent area
allocation, without the circle-slice certificate.

\section{An explicit connection in every degree of a Blaschke-power family}

The family comes from the round-five corpus. The new point is a uniform
finite theorem with no unspecified starting degree.
For $0<b<1$, put
\[
 A(z)=z(z+b),\quad D(z)=1+bz,\quad
 B(z)=\frac{A(z)}{D(z)},\quad F_{N,b}(z)=A(z)^N-D(z)^N.
\]
The finite Blaschke-product interpretation of $B$ is standard
\cite{blaschke}. The argument below also proves the required root geometry
by direct algebra.

\begin{theorem}[quadratic-sector connection]\label{thm:sector}
Let $N\ge3$ and $0<b\le1/(10N)$, and define
\[
 \lambda=Nb,\qquad m=2N-1,\qquad
 t=\left(\frac{2\lambda}{5}\right)^{1/m},\qquad q=\frac\pi N.
\]
The monic polynomial $F_{N,b}$ has $2N$ simple roots on the unit circle.
Two of them are joined by an explicitly parametrised, piecewise smooth
curve $\gamma$ satisfying
\begin{align}
 \max_\gamma |F_{N,b}|&\le \exp(-\lambda t/25)<1,\label{eq:collar}\\
 \length(\gamma)&\le\mathcal L_{N,b}(t)\nonumber\\
 &:=2(1-t)+qt
    +b(1-t^2+qt^2)
    +3b^2\left(2(t^{-1}-1)+q/t\right)\label{eq:exactlength}\\
 &\le2-\frac{689}{1050}t<2.\label{eq:uniformlength}
\end{align}
For every $0<r<1$, the monic polynomial
$G(z)=r^{2N}F_{N,b}(z/r)$ has simple roots strictly inside the unit disc,
and $r\gamma$ has length $r\length(\gamma)$ and level at most
$r^{2N}\exp(-\lambda t/25)$.
\end{theorem}

\subsection{Root placement and a low-degree coordinate}

For $z=x+iy$ and $R=|z|$,
\begin{equation}\label{eq:normdefect}
 |A(z)|^2-|D(z)|^2
 =(R^2-1)(R^2+1+2bx).
\end{equation}
The second factor equals $(x+b)^2+y^2+1-b^2>0$.
At a root of $F_{N,b}$, $|A|=|D|$, so $|z|=1$.
Moreover $A,D$ have no common zero for $b<1$.
If this root were multiple, $F=F'=0$ would imply
\[
 DA'-AD'=bz^2+2z+b=0.
\]
On $|z|=1$, however, $|2z|=2>2b\ge|b(z^2+1)|$.
Thus all $2N$ roots are simple and lie on the circle.

Let
\[
 \Sigma=\{p:t\le|p|\le1,\ 0\le\arg p\le\pi/N\}.
\]
First observe that $t>10b$. Indeed,
\[
 \frac{(10b)^{2N-1}}{(2/5)Nb}
 \le\frac{25}{N^{2N-1}}\le\frac{25}{3^5}<1.
\]
On a neighbourhood of $\Sigma$, define
\begin{equation}\label{eq:H}
 e(p)=\frac{b^2}{4}(p^{-1}-p)^2,\quad
 s(p)=\sqrt{1+e(p)},\quad
 H(p)=\frac b2(p^2-1)-ps(p),
\end{equation}
using the principal square root. If $u=|p|\in[t,1]$, then
\begin{equation}\label{eq:small-e}
 |e(p)|\le b^2/u^2\le1/100,\quad
 \Rea s(p)>0,\quad |s(p)|\ge99/100,\quad
 |s(p)-1|\le|e(p)|.
\end{equation}
The last inequality follows from $(s-1)(s+1)=e$ and $|s+1|\ge1$.
All the formulae are analytic on a neighbourhood of the closed sector.

The quadratic formula gives the exact identities
\begin{equation}\label{eq:pullback}
 H^2+b(1-p^2)H-p^2=0,\qquad
 A(H)=p^2D(H),\qquad
 F_{N,b}(H)=D(H)^N(p^{2N}-1).
\end{equation}
Equation~\eqref{eq:normdefect} and $|p|\le1$ imply $|H(p)|\le1$.
Consequently $D(H)\ne0$. The points $H(1)=-1$ and
$H(e^{i\pi/N})$ are roots of $F_{N,b}$, and they are distinct because
their values of $A/D=p^2$ differ.

This is the change of variable responsible for the theorem. The inverse
belongs to the quadratic map $B$, rather than to the high-degree
polynomial $F_{N,b}$. Crowding of critical values of $F_{N,b}$ therefore does
not limit this inverse neighbourhood.

\subsection{The amplitude deficit pays for the turn}

Write $p=ue^{i\theta}\in\Sigma$. Since $\theta\le\pi/3$,
$\Rea\{(b/2)(p^2-1)\}\le0$, and~\eqref{eq:small-e} gives
\[
 \Rea H(p)\le-u\cos\theta+b^2/u
 \le-\frac12u+\frac1{100}u=-\frac{49}{100}u.
\]
Together with $|H|\le1$ and $b\le u/10$, this proves
\begin{align*}
 |D(H)|^2
 &\le1-\frac{49}{50}bu+b^2
 \le1-\frac{22}{25}bu,\\
 |D(H)|^N&\le\exp\left(-\frac{11}{25}\lambda u\right).
\end{align*}

Take the path in the $p$-plane consisting of
\[
 1\longrightarrow t\longrightarrow te^{i\pi/N}
 \longrightarrow e^{i\pi/N},
\]
with two radial segments and the shorter circular arc. Define $\gamma$ as
its image under $H$.
On either radial piece, $p^{2N}=u^{2N}\in[0,1]$, so
$|p^{2N}-1|\le1$. On the circular piece,
\[
 |p^{2N}-1|\le1+t^{2N}
 =1+\frac25\lambda t.
\]
Hence, throughout the path,
\[
 |F_{N,b}(H(p))|
 \le e^{-11\lambda t/25}(1+2\lambda t/5)
 \le e^{-\lambda t/25}<1.
\]
This proves~\eqref{eq:collar}. Taking absolute values of $A^N$ and $D^N$
separately would destroy the radial cancellation in~\eqref{eq:pullback}.
The turn becomes affordable only after the path has accumulated a strict
amplitude deficit.

\subsection{Integrate the actual derivative bound}

Let $E(p)=H(p)+p$. From~\eqref{eq:H},
\[
 E'=bp+1-s-\frac{pe'}{2s},\qquad
 |pe'|\le2b^2/u^2.
\]
Thus
\begin{equation}\label{eq:derivative}
 |E'(p)|\le bu+3b^2/u^2,
 \qquad |H'(p)|\le1+bu+3b^2/u^2.
\end{equation}
The coefficient $3$ is deliberately nonoptimal:
$1+100/99<3$. Integrating~\eqref{eq:derivative} on the two radial pieces
and the circular piece yields~\eqref{eq:exactlength}.

Since $q\le22/21$ and $b/t\le1/10$, the first two terms are at most
$2-(20/21)t$. The remaining terms are bounded above by
\[
 b(1+q)+\frac{3b^2}{t}(2+q)
 \le\left(\frac{43}{210}+\frac{96}{1050}\right)t
 =\frac{311}{1050}t.
\]
Subtracting gives~\eqref{eq:uniformlength} and completes the proof of
Theorem~\ref{thm:sector}. The contraction statement is direct.

\begin{corollary}[the previous asymptotic regime, with an explicit threshold]
For fixed $0<\lambda\le1/10$ and $b=\lambda/N$, every $N\ge3$ is covered,
and the refined estimate gives
\[
 \length(\gamma)
 \le\frac{\log(5/(2\lambda))+\pi+o(1)}{N}.
\]
\end{corollary}
\begin{proof}
Here $t=\exp[-\log(5/(2\lambda))/(2N-1)]$. In
\eqref{eq:exactlength}, the $b$-term is $O(N^{-2})$ and the $b^2$-term is
$O(N^{-3})$, while the first two terms have the displayed expansion.
\end{proof}

\section{A finite high-critical witness with a path shorter than one half}

The round-five certificate concerns
\begin{equation}\label{eq:G24}
 N=12,\quad b=1/120,\quad r=999999999/10^9,\quad
 G(z)=r^{24}F_{12,1/120}(z/r).
\end{equation}
All its roots are simple and have modulus $r<1$.
The existing exact rational certificate, replayed in this review, encloses
all 23 critical points in disjoint discs and proves
\begin{align}
 |G(c)+1|&<1/12,\label{eq:oldvalues}\\
 \frac1{23}\sum_{G'(c)=0}|G(c)|^{8/23}
 &>\frac{46001266667}{46000000000}
 >1+\frac1{40000}.\label{eq:oldmoment}
\end{align}
These are pre-existing spectral claims, not new deductions from a numerical
root finder. Equation~\eqref{eq:oldvalues} implies $\mu(G)>11/12$ and, for
any pair of critical values,
\[
 |v_i/v_j-1|<2/11.
\]
Thus $|v_i/v_j-a|<13/11<4/3$ for every $a\in[0,1]$.
No selected critical value passes the displayed radius-$4/3$ test.

\begin{corollary}[new finite metric consequence]\label{cor:24}
Two roots of the very same polynomial~\eqref{eq:G24} have an explicit curve
$\gamma_G$ with
\[
 \length(\gamma_G)
 <\frac{58672129}{117936000}<\frac12,
 \qquad
 \max_{\gamma_G}|G|\le\frac{3750}{3763}<1.
\]
\end{corollary}
\begin{proof}
Now $\lambda=1/10$ and $t=(1/25)^{1/23}>13/15$, since
$(13/15)^{23}<1/25$. In~\eqref{eq:exactlength}, for these fixed $N,b$,
every term in
\[
 2-(2-q)t+b[1-(1-q)t^2]+3b^2[(2+q)/t-2]
\]
is decreasing in $t$, and the expression is increasing in $q$.
Use $q<11/42$ and $t>13/15$. The exact resulting upper bound is
\[
 2-\left(2-\frac{11}{42}\right)\frac{13}{15}
 +\frac1{120}\left[1-\left(1-\frac{11}{42}\right)
                 \left(\frac{13}{15}\right)^2\right]
 +\frac3{120^2}\left[2\left(\frac{15}{13}-1\right)
                    +\frac{11}{42}\frac{15}{13}\right]
 =\frac{58672129}{117936000}<\frac12.
\]
Furthermore $\lambda t/25>13/3750$ and $e^{-x}\le(1+x)^{-1}$ for $x\ge0$.
Equation~\eqref{eq:collar} gives the stated level. Contraction by $r$ only
improves both inequalities.
\end{proof}

This is a quantitative failure of two sufficient tests to characterise
metric difficulty. It is not a counterexample to the parent problem, nor
to the numerator-four moment theorem. It also does not show that no theorem
using more complete critical-spectrum information can control paths.
Its precise assertion is that the tested high-critical, nonseparated regime
can contain a polynomial with a connection shorter than one half.

\section{Weighted radial barriers and an exact cusp order}

For a monic polynomial with roots in the closed unit disc, let
\[
 \Lambda(f)=\inf\{\length(\gamma):\gamma\subset K_1(f)
     \text{ joins two different listed root occurrences}\}.
\]
For $z^n-1$ the punctured sublevel has separate root sectors. Every connection
between distinct roots passes through zero, so $\Lambda(z^n-1)=2$.

\subsection{A reusable barrier lemma}

\begin{lemma}[weighted perturbation barrier]\label{lem:barrier}
Let $n>j\ge1$, $0<\eta<1$, and $P(z)=z^n-1+E(z)$.
Suppose all roots are simple, lie on the unit circle, and there is exactly
one root in each sector between the rays
$\arg z=(2k+1)\pi/n$. Assume
\[
 |E(z)|\le\eta |z|^j\quad(|z|\le1),\qquad
 |E(z)|\le\eta |z|^{n-1}\quad(|z|\ge1).
\]
Then
\[
 \Lambda(P)\ge2\left(1-\eta^{1/(n-j)}\right).
\]
\end{lemma}
\begin{proof}
On a separating ray, $z^n=-|z|^n$, so
$|z^n-1|=1+|z|^n$. For $\eta^{1/(n-j)}<|z|\le1$,
\[
 |P(z)|\ge1+|z|^n-\eta |z|^j>1.
\]
For $|z|\ge1$ the same conclusion follows from
$|z|^n-\eta |z|^{n-1}>0$. The rays outside radius
$R=\eta^{1/(n-j)}$ are therefore barriers extending to infinity.
A path between roots in different sectors must enter $|z|\le R$.
Its portions before and after that first inward excursion have total length
at least $2(1-R)$.
\end{proof}

The vanishing order of $E$ matters. A uniform bound $|E|\le\eta$ would
retain only the exponent $1/n$. Retaining the zero of the perturbation at
the origin gives the exponent $1/(n-1)$ used next.

\subsection{Apply the lemma to the sector family}

Put $n=2N$ and $E_b=F_{N,b}-(z^{2N}-1)$.
The coefficient of $z^N$ cancels exactly, and all the other nonzero
coefficients in the upper and lower halves have opposite signs. Hence
\begin{equation}\label{eq:delta}
 \delta_b:=\|E_b\|_{\mathrm{coeff},1}
 =2\bigl((1+b)^N-1-b^N\bigr),\qquad
 2\lambda\le\delta_b\le\frac{20}{9}\lambda.
\end{equation}
For the upper bound use $\binom Nk\le N^k$, then sum the geometric series:
$\sum_{k\ge1}\lambda^k\le\lambda/(1-\lambda)$ with $\lambda\le1/10$.
Every monomial of $E_b$ has degree between $1$ and $2N-1$, so
\[
 |E_b(z)|\le\delta_b|z|\quad(|z|\le1),\qquad
 |E_b(z)|\le\delta_b|z|^{2N-1}\quad(|z|\ge1).
\]
These bounds also show that no zero crosses a separating ray as $b$ varies
from zero to its final value. All roots remain simple and on the circle;
at $b=0$ there is exactly one in each sector. Continuity of the simple roots
therefore gives the sector-count hypothesis for all permitted $b$.

\begin{theorem}[two-sided fixed-degree cusp]\label{thm:cusp}
For the family of Theorem~\ref{thm:sector}, with $n=2N$, $m=n-1$ and
$\delta_b$ as in~\eqref{eq:delta},
\begin{equation}\label{eq:cusp}
 \boxed{
 \frac{689}{1050}\left(\frac{9\delta_b}{50}\right)^{1/m}
 \le 2-\Lambda(F_{N,b})
 \le 2\delta_b^{1/m}.}
\end{equation}
In particular, for every fixed even $n\ge6$,
$2-\Lambda(F_{N,b})\asymp\delta_b^{1/(n-1)}$ as $b\downarrow0$.
No local H\"older estimate at $z^n-1$ with exponent greater than
$1/(n-1)$ holds for $\Lambda$ on the full closed-root-disc coefficient class.
\end{theorem}
\begin{proof}
The lower bound on $\Lambda$ follows from Lemma~\ref{lem:barrier} with
$\eta=\delta_b<1$ and $j=1$, so $2-\Lambda\le2\delta_b^{1/m}$.
The constructed curve gives
$2-\Lambda\ge(689/1050)(2\lambda/5)^{1/m}$.
Since $\lambda\ge9\delta_b/20$, this implies the other half of
\eqref{eq:cusp}. If a local estimate
$|\Lambda(f)-\Lambda(z^n-1)|\le C\|f-(z^n-1)\|_{\mathrm{coeff},1}^{\alpha}$
held with $\alpha>1/(n-1)$, substitution of this family and division by
$\delta_b^{1/(n-1)}$ would contradict the positive lower constant.
\end{proof}

\begin{remark}[relation to the supplied numerical observation]
\path{StraightSpokeHubCriterionLab.md}, \S6d, measures a cusp exponent
$1/(n-1)$ for a normalised canonical-branch statistic in several degrees.
Theorem~\ref{thm:cusp} proves that order for the all-curve functional
$\Lambda$, on an explicit family in every even degree $n\ge6$.
It does not prove that a canonical arc minimises $\Lambda$, or establish a
H\"older upper modulus for arbitrary perturbations of the binomial.
\end{remark}

The constants can be tightened. In fact the weighted barrier above uses
$\eta=\delta_b$, rather than the rougher $(20/9)\lambda$.
Neither a leading coefficient for the cusp nor an exact minimising curve is
claimed here. The two-sided power law is already enough to exclude a
Lipschitz or Hessian-scale argument at this singular configuration.

\subsection{A natural rescaled geometry}

Fix $n=2N$ and put $\varepsilon=\lambda^{1/(n-1)}$. Direct expansion gives,
uniformly for $w$ in compact sets,
\begin{equation}\label{eq:blowup}
 \frac{F_{N,b}(\varepsilon w)+1}{\varepsilon^n}
 \longrightarrow w^n-w.
\end{equation}
The first term from $A^N-z^n$ is $\lambda z^{n-1}$, of size
$\varepsilon^{2n-2}$ after substitution; the first term from $D^N-1$ is
$\lambda z=\varepsilon^nw$. The remaining terms are of higher order.
Thus the formal limiting admissible region is
\[
 \Omega_n=\{w:\Rea(w^n-w)\ge0\}.
\]
Indeed, writing the left side of~\eqref{eq:blowup} as $Q_b(w)$, the exact
scaled containment condition is
\[
 \Rea Q_b(w)\ge\frac{\varepsilon^n}{2}|Q_b(w)|^2.
\]
For even $n$, no critical point of $w^n-w$ lies on the zero real-part level,
so the finite-plane boundary is regular. Its critical points satisfy
$nw^{n-1}=1$, and a purely imaginary such point cannot occur because $n-1$
is odd. This does not by itself prove convergence of the renormalised
intrinsic distances, whose endpoints recede to infinity.

A precise follow-up is whether
\[
 \lim_{b\downarrow0}
 \frac{2-\Lambda(F_{N,b})}{(Nb)^{1/(n-1)}}
\]
exists, and whether it is the renormalised distance between ends of
$\Omega_n$. Equation~\eqref{eq:cusp} supplies positive finite bounds;
identification of this variational limit remains unproved here.

For degree tending to infinity, the same explicit upper path also recovers
the registered exponential degeneration. With $\lambda_N=e^{-\alpha N}$,
$t_N\to e^{-\alpha/2}$ and the error terms in~\eqref{eq:exactlength} vanish.
The weighted barriers give the matching lower limit. Consequently
$\Lambda(F_{N,\lambda_N/N})\to2(1-e^{-\alpha/2})$, consistent with the
round-five theorem. The new estimate supplies a uniform finite precursor
and the stronger fixed-degree cusp assertion.

\section{The mechanism hidden behind the live $13/25$ theorem}

The most reusable registered geometric idea is the passage from angular
projection to common circle slices. The mechanism belongs in the short note
before the implementation details of its exact certificate.

Under the no-short-connection assumption, a conformal normalisation of an
ancestor component produces points $b_1,\ldots,b_k$ in the hyperbolic disc.
With $a=\Area(C_t)/\pi$, the source gives pairwise separation
\[
 d(b_i,b_j)\ge D=4\operatorname{artanh}\sqrt{\tanh(1/a)},
 \qquad \cosh(D/2)=e^{2/a},
\]
a lower radial bound $d_j=d(0,b_j)\ge d_*(a)>0$, and
\[
 \sum_j\ell(d_j)\ge x=\log(t/\mu),\qquad
 \ell(d)=-\log\tanh(d/2).
\]
The exact formula for $d_*$, the point supplying the normalisation, and the
comparison differential inequality remain in the complete registered proof.

\begin{lemma}[common-slice dual bound]\label{lem:slice}
Let $b_1,\ldots,b_k$ be pairwise hyperbolically $D$-separated points with
$d_j\ge d_*>0$. Define
\[
 w(d,r)=\arccos\clip{\frac{\cosh d\cosh r-\cosh(D/2)}
                              {\sinh d\sinh r}}.
\]
Then $\sum_jw(d_j,r)\le\pi$ for every $r>0$. If radii $r_i>0$ and
weights $\sigma_i\ge0$ satisfy
\[
 \ell(d)\le U+\sum_i\sigma_iw(d,r_i)\quad(d\ge d_*),
 \qquad U>0,
\]
then $\sum_j\ell(d_j)\le kU+\pi\sum_i\sigma_i$.
\end{lemma}
\begin{proof}
The open hyperbolic balls of radius $D/2$ about the $b_j$ are disjoint.
Their intersections with a fixed hyperbolic circle about zero are disjoint
arcs. The hyperbolic cosine law gives angular lengths $2w(d_j,r)$, including
empty and full-circle cases. Sum their lengths, then sum the proposed
majorant over $j$ and use the nonnegative weights.
\end{proof}

This is an explicit geometric dual certificate. Any nonnegative weights
work once their majorant is proved. Optimisation proposes weights; it does
not confer proof authority. The reported $13/25$ certificate has been
replayed in full mode in this review: 126 duals, with
$X_{\rm cert}=635762889599/10^{12}$ and threshold $13/25$.

The failed predecessor uses only consecutive angular gaps. It allows an
alternating radial profile $d_*,d_*+D,d_*,d_*+D,\ldots$.
Every mixed neighbouring pair has radial difference $D$ and therefore costs
zero forced angle. The scalar budget admits arbitrarily many nearby roots
separated in its cyclic list by distant decoys. This is a failure of the
relaxation; the profile need not come from an actual hyperbolic packing.
Common slices retain the all-pair disjointness that the relaxation discarded.

The construction in Section~2 has a related logical feature. It retains the
quadratic pullback constraint before bounding the high power. In both
arguments, a simplification made too early destroys exactly the relation
needed to join the estimates.

\section{Revision architecture and the surviving obstruction}

\subsection{What is already improved}

The supplied live note has 24 pages. Its trinomial theorem and complete
proof already finish on page one. Its Abel control-polygon paragraph already
explains the unrestricted coefficient. These changes should be retained,
rather than recommended again as though the prior edits had not landed.
The exact sextic remains a useful boundary immediately afterwards.

The applicable patch preserves the title, the first theorem and its complete
proof byte for byte, and the public commit pin. It retains the $13/25$,
$71/10$ and radius-$4/3$ results. It inserts only the registered scaling
consequence, not the new sector or cusp theorems pending adjudication.

The main changes are mathematical rather than typographical. The strongest
general bound is made visible. The low-critical proof is explained through
circle slices and a finite dual majorant. Numerically subsumed arity/capacity
conclusions move to the long record with their component-sensitive meaning
intact. Repeated verification inventories are replaced by a compact guide
beginning with the actual first theorem. The complete removed text is
archived alongside the patch.

\subsection{Two local corrections}

The parent separation corollary must state $0<|v|<1$, since it divides by
$v$. The old wording gives only $|v|<1$ and is not sufficient in the
presence of multiple roots.

The one-root perimeter discussion currently invokes $z^N-z$ at level one.
Its critical values all have modulus below one, so it cannot furnish the
claimed one-root component there. The correct family in the underlying
source is $z^N-az$ with fixed $a>1$, then $N\to\infty$ and $a\downarrow1$.
For any $r<1/a$, the central component eventually contains $|z|\le r$.
A circle of radius between $1/a$ and one isolates exactly one root by
Rouch\'e. This proves the limiting perimeter lower bound $2\pi/a$;
letting $a\downarrow1$ gives the necessary universal constant $2\pi$.
It proves no corresponding upper bound.

\subsection{A structural map of insufficient information}

\begin{center}\small
\begin{tabular}{p{.24\linewidth}p{.32\linewidth}p{.34\linewidth}}
\toprule
Information retained & Exact failure & Information a successor must retain\\
\midrule
Partial sums only implicitly & Sextic leaves a prescribed spoke & Explicit convex control of the selected spoke, or another curve\\
Consecutive angular gaps & Alternating radii evade the relaxed budget & Disjointness on common slices\\
Initial merge arity & Later mergers erase a capacity gap & Component geometry at the actual working level\\
Angular mean length & Attachable directions can occupy a shrinking window & Attachment-conditioned estimates or freely chosen attachments\\
Two scalar critical-value tests & Degree-24 witness fails both but has length below $1/2$ & A coordinate or selection retaining geometric cost\\
A Hessian-scale perturbation model & Two-sided $\delta^{1/(n-1)}$ cusp & Singular scaling and the relevant intrinsic geometry\\
\bottomrule
\end{tabular}
\end{center}

The exact all-curve question remains $\Lambda(f)\le2$ on the closed-root-disc
class. The registered lower-semicontinuity theorem passes a bound from a
dense generic subclass to the whole fixed-degree class. Scaling an enclosing
disc of radius $R<1$ then gives a strict length below two and strict
unit-level containment. This argument does not require a uniform forward
transport of an already selected curve.

The generic slit-sheet construction is already an ordinary theorem. A
sufficient remaining route is
\[
 \min_{\substack{f'(c)=0\\|f(c)|\le1}}L_f(c)\le2,
\]
where $L_f(c)$ is the admissible canonical root-to-root descent length.
This minimum bounds $\Lambda$ from above; equality is not established.
The new sector curve is not forced to be a Newton descent of $F_{N,b}$.
It illustrates why the minimisation should remain open to other competitors.

The source route's assertions that FP$_m$ for $m\ge5$ remains open and
that generic slit construction is still missing are obsolete mathematical
frontiers. The relevant ordinary proof records are, respectively:
\begin{quote}\small
\path{FreePointQuadraticAllDegrees.md};\\
\path{AttachmentAwareReeb.md};\\
\path{GenericSufficiencyClosure.md}.
\end{quote}
Formalising their full analytic content is a separate task.

\section{Further questions, antecedents and proof assurance}

The most concrete new question is the rescaled variational limit in
Section~4. A more general construction question is to identify conditions on
an analytic coordinate $H$ and amplitude $a$ under which
\[
 f(H(p))=a(p)(p^m-1)
\]
admits a short radial--circular--radial connection. The sufficient ingredients
in this proof are endpoint phases, an amplitude deficit on a common sector,
and an integrable derivative bound. None of these is automatically available
for a general polynomial.

For the low-critical route, the useful optimisation is now a better
majorant on common slices or a sharper comparison using that same geometry.
Returning to angular shadows would reintroduce the proved failure.
For the remaining high-critical class, a theorem must select a curve together
with its admissibility and length; separate marginal statements cannot be
combined without a compatibility argument.

The connected \#249 source gives another example. Its file
\begin{quote}\small\path{TotientActualLcmOrbitSeparation.lean}\end{quote}
ends with an irrationality theorem whose actual-orbit separation supply is
an explicit hypothesis. A
favourable abstract residue and a bound on the actual orbit are not by
themselves a jointly attained margin. This is a shared proof-design issue,
not a theorem transferring the present geometry to the totient problem.
No sibling source is edited in this return.

Tao's partial-progress discussion stresses the precise point where an
unsuccessful technique stops working \cite{taopartial}; his
results-to-effort guidance supports including cheap consequences of live
results and retaining illuminating examples \cite{taoeffort}.
Here that means showing Proposition~\ref{thm:scale}, the circle-slice
mechanism, and a finite obstruction that is also explicitly solved.
A list of model or verification counts does not convey these arguments.

Pendyala's degree-four theorem directly concerns the parent root-pair
question \cite{pendyala}. The Blaschke-product root geometry is classical
\cite{blaschke}. The polynomial family and its asymptotic spectral and
metric behaviour are already in the round-five corpus.
The uniform finite sector proof and fixed-degree all-curve cusp bounds above
are the additions of this review. The literature search did not establish
their novelty. Neither a failed search nor a formal proof would establish
priority.

Palomar's current requirements distinguish an independently stated Challenge,
a proved Solution and an immutable source revision; they do not turn a
finite scalar component into a complete analytic theorem \cite{palomar}.
The provided Lean module contains proposed algebraic and inequality
components and no admitted theorem. It is uncompiled here. It does not claim
to formalise analytic square roots, the integrated path, sector topology or
intrinsic-distance minimisation.

\paragraph{Executed checks.}
The full angular-budget checker passed with 126 duals and threshold $13/25$.
The degree-24 rational critical-value certificate passed again.
The new rational checker verifies the radius, length and collar constants,
the homogeneity comparison and finite coefficient regressions.
An optional SymPy script checks three polynomial identities exactly.
These checks are accompanied by their outputs and reproduction commands.
The ordinary geometry is proved in this document rather than certified by
sampling. No Lean compilation, Comparator run or independent human review is
claimed.

\begin{thebibliography}{9}
\bibitem{blaschke} S.~R. Garcia, J.~Mashreghi and W.~T. Ross,
\emph{Finite Blaschke products: a survey}, arXiv:1512.05444v2 (2016);
Harmonic Analysis, Function Theory, Operator Theory, and their Applications,
Theta Series in Advanced Mathematics (2017), 133--158.
\url{https://arxiv.org/abs/1512.05444}.
\bibitem{pendyala} V.~S. Pendyala, \emph{A Degree-Four Lemniscate Path Theorem},
arXiv:2606.24875v1 (2026), Theorem~1.
\url{https://arxiv.org/abs/2606.24875}.
\bibitem{taopartial} T. Tao, \emph{On the importance of partial progress},
originally posted 17 July 2012, accessed 7 September 2026.
\url{https://terrytao.wordpress.com/career-advice/on-the-importance-of-partial-progress/}.
\bibitem{taoeffort} T. Tao, \emph{Maximising the results-to-effort ratio},
accessed 7 September 2026.
\url{https://terrytao.wordpress.com/advice-on-writing-papers/maximising-the-results-to-effort-ratio/}.
\bibitem{palomar} Palomar, \emph{How to submit}, accessed 7 September 2026.
\url{https://palomar-registry.org/how-to-submit}.
\bibitem{corpus} W. Cook and the Plectis corpus, round-six supplied packet,
\path{01_short_note.tex}, \path{03_research_packet.json}, and the accompanying
\#1041 source slice. Public browse revision
\texttt{32545d77c4b5cfd72ebd36c8dcd2418bf7cc4d31}; the supplied live source
has precedence over the public snapshot. All ordinary-source filenames in
this return are relative to \path{ErdosProblems/Erdos1041/} unless stated.
\end{thebibliography}
\end{document}
